\documentclass[12pt,a4paper,reqno]{amsart}

% language
\usepackage[greek,english]{babel}
\usepackage[utf8]{inputenc}
\usepackage{alphabeta}

% change default names to greek
\addto\captionsenglish{
    \renewcommand{\contentsname}{Περιεχόμενα}
    \renewcommand{\refname}{Βιβλιογραφία}
    \renewcommand{\datename}{Ημερομηνία:}
    \renewcommand{\urladdrname}{Ιστοσελίδα}
}

% math 
\usepackage{amsmath,amsthm,amssymb,amscd}

% font
\usepackage[cal=euler]{mathalfa}
\usepackage{libertinus-type1}
% \usepackage{txfonts} % for upright greek letters
\usepackage{bm} % for bold symbols
\usepackage{bbm} % for the simply-looking bb symbols

% miscellaneous 
\usepackage{changepage} %for indenting environments
\usepackage{csquotes} % example: \textcquote{}
\usepackage{blkarray}
\setcounter{MaxMatrixCols}{20} % default for pmatrix is 10!!
\usepackage{ytableau}
\usepackage{array} %needed to increase the vertical length in a tabular

% drawing
\usepackage{tikz,tikz-cd}
\usetikzlibrary{shapes.misc, patterns, matrix, calc, intersections,positioning,arrows.meta}
\usepackage{graphics,graphicx}
\usepackage{float} % provides enhanced control and customization options for floating objects such as figures and tables

% colors
\usepackage{xcolor}
\definecolor{darkcandyapplered}{rgb}{0.64, 0.0, 0.0}
\definecolor{midnightblue}{rgb}{0.1, 0.1, 0.44}
\definecolor{mylightblue}{HTML}{336699}
\definecolor{burntorange}{rgb}{0.8, 0.33, 0.0}
\definecolor{iceberg}{rgb}{0.44, 0.65, 0.82}
\definecolor{applegreen}{rgb}{0.55, 0.71, 0.0}
\definecolor{canaryyellow}{rgb}{1.0, 0.94, 0.0}
\definecolor{lightgray}{rgb}{0.83, 0.83, 0.83}

% hrefs
\usepackage{hyperref}
\usepackage[noabbrev,capitalize]{cleveref}
\hypersetup{
    pdftoolbar=true,        
    pdfmenubar=true,        
    pdffitwindow=false,     
    pdfstartview={FitH},  % fits the width of the page to the window
    pdftitle={},
    pdfauthor={},
    pdfsubject={},
    pdfkeywords={},
    pdfnewwindow=true,  % links in new window
    colorlinks=true,  % false: boxed links; true: colored links
    linkcolor=darkcandyapplered,   % color of internal links
    citecolor=midnightblue,  % color of links to bibliography
    urlcolor=cyan,  % color of external links
    linktocpage=true  % changes the links from the section body to the page number
    }

% geometry
\textwidth=16cm 
\textheight=21cm 
\hoffset=-55pt 
\footskip=25pt

% thm envs (you might need to change the path)
\theoremstyle{definition}
\newtheorem{exercise}{Άσκηση}
\newtheorem{solution}{Λύση}
\newtheorem*{remark}{Παρατήρηση}
\newtheorem*{example}{Παράδειγμα}

% math ops 
% fields
\newcommand{\NN}{\mathbbmss{N}} 
\newcommand{\ZZ}{\mathbbmss{Z}} 
\newcommand{\QQ}{\mathbbmss{Q}} 
\newcommand{\RR}{\mathbbmss{R}} 
\newcommand{\CC}{\mathbbmss{C}} 
\newcommand{\KK}{\mathbbmss{K}} 
\newcommand{\FF}{\mathbbmss{F}} 
\newcommand{\HH}{\mathbbmss{H}} 

% symmetric group
\newcommand{\fS}{\mathfrak{S}}  
\newcommand{\fp}{\mathfrak{p}}  
\newcommand{\fm}{\mathfrak{m}}  

% calligraphic 
\newcommand{\aA}{\mathcal{A}} 
\newcommand{\bB}{\mathcal{B}}
\newcommand{\cC}{\mathcal{C}}
\newcommand{\dD}{\mathcal{D}}
\newcommand{\eE}{\mathcal{E}}
\newcommand{\fF}{\mathcal{F}}
\newcommand{\hH}{\mathcal{H}}
\newcommand{\iI}{\mathcal{I}}
\newcommand{\lL}{\mathcal{L}}
\newcommand{\nN}{\mathcal{N}}
\newcommand{\oO}{\mathcal{O}}
\newcommand{\pP}{\mathcal{P}}
\newcommand{\sS}{\mathcal{S}}
\newcommand{\mM}{\mathcal{M}}
\newcommand{\uU}{\mathcal{U}}

% bold
\newcommand{\bfa}{\mathbf{a}}
\newcommand{\bfe}{\mathbf{e}}
\newcommand{\bfF}{\pmb{F}}
\newcommand{\bfR}{\pmb{R}}
\newcommand{\bfv}{\mathbf{v}}
%\newcommand{\bfx}{\bm{x}}
%\newcommand{\bfx}{\mathbf{x}} 
\newcommand{\bfx}{\pmb{x}}
\newcommand{\bfX}{\pmb{X}}
\newcommand{\bfy}{\pmb{y}}
\newcommand{\bfz}{\pmb{z}}

% roman
\newcommand{\rmA}{\mathrm{A}}
\newcommand{\rmB}{\mathrm{B}}
\newcommand{\rmC}{\mathrm{C}}
\newcommand{\rmD}{\mathrm{D}} 
\newcommand{\rmH}{\mathrm{H}} 
\newcommand{\rmh}{\mathrm{h}} 
\newcommand{\rmI}{\mathrm{I}} 
\newcommand{\rmJ}{\mathrm{J}} 
\newcommand{\rmK}{\mathrm{K}}
\newcommand{\rmM}{\mathrm{M}}
\newcommand{\rmP}{\mathrm{P}}  
\newcommand{\rmp}{\mathrm{p}}  
\newcommand{\rmQ}{\mathrm{Q}}  
\newcommand{\rmR}{\mathrm{R}}
\newcommand{\rmS}{\mathrm{S}}
\newcommand{\rmT}{\mathrm{T}}
\newcommand{\rmU}{\mathrm{U}}
\newcommand{\rmV}{\mathrm{V}}
\newcommand{\rmY}{\mathrm{Y}}
\newcommand{\rmZ}{\mathrm{Z}}
\newcommand{\rmz}{\mathrm{z}}

% arrows and symbols 
\renewcommand{\to}{\rightarrow}
\newcommand{\toto}{\longrightarrow}
\newcommand{\mapstoto}{\longmapsto}
\newcommand{\then}{\Rightarrow}
\newcommand{\IFF}{\Leftrightarrow}
\newcommand{\tl}{\tilde}
\newcommand{\wtl}{\widetilde}
\newcommand{\ol}{\overline}
\newcommand{\ul}{\underline}
\newcommand{\oldemptyset}{\emptyset}
\renewcommand{\emptyset}{\varnothing}
\DeclareMathSymbol{\Arg}{\mathbin}{AMSa}{"39} % for arguments 
\newcommand{\onto}{\ensuremath{\twoheadrightarrow}}
\newcommand{\tle}{\trianglelefteq}
\newcommand{\tge}{\trianglerighteq}

% custom ops
\newcommand{\rmKer}{\mathrm{Ker}}
\newcommand{\rmIm}{\mathrm{Im}}
\newcommand{\lcm}{\mathrm{lcm}}
\newcommand{\GL}{\mathrm{GL}}
\newcommand{\ev}{\mathrm{ev}}
\newcommand{\End}{\mathrm{End}}
\newcommand{\rmid}{\mathrm{id}}
\newcommand{\rmspan}{\mathrm{span}}
\newcommand{\Spec}{\mathrm{Spec}}
\newcommand{\mSpec}{\mathrm{mSpec}}
\newcommand{\rad}{\mathrm{rad}}
\newcommand{\Rad}{\mathrm{Rad}}
\newcommand{\Jac}{\mathrm{Jac}}
\newcommand{\tor}{\mathrm{tor}}
\newcommand{\Ann}{\mathrm{Ann}}
\newcommand{\Hom}{\mathrm{Hom}}
\newcommand{\Hilb}{\mathrm{Hilb}}

% absolute value symbol
\usepackage{mathtools} 
\DeclarePairedDelimiter\abs{\lvert}{\rvert}%
\DeclarePairedDelimiter\norm{\lVert}{\rVert}%
\makeatletter
\let\oldabs\abs
\def\abs{\@ifstar{\oldabs}{\oldabs*}}

% tensor symbol
\newcommand{\tensor}[1]{%
  \mathbin{\mathop{\otimes}\limits_{#1}}%
}

% permutation cycle notation
\ExplSyntaxOn
\NewDocumentCommand{\cycle}{ O{\;} m }
 {
  (
  \alec_cycle:nn { #1 } { #2 }
  )
 }

\seq_new:N \l_alec_cycle_seq
\cs_new_protected:Npn \alec_cycle:nn #1 #2
 {
  \seq_set_split:Nnn \l_alec_cycle_seq { , } { #2 }
  \seq_use:Nn \l_alec_cycle_seq { #1 }
 }
\ExplSyntaxOff

% setminus symbol
\newcommand{\mysetminusD}{\hbox{\tikz{\draw[line width=0.6pt,line cap=round] (3pt,0) -- (0,6pt);}}}
\newcommand{\mysetminusT}{\mysetminusD}
\newcommand{\mysetminusS}{\hbox{\tikz{\draw[line width=0.45pt,line cap=round] (2pt,0) -- (0,4pt);}}}
\newcommand{\mysetminusSS}{\hbox{\tikz{\draw[line width=0.4pt,line cap=round] (1.5pt,0) -- (0,3pt);}}}
\newcommand{\sm}{\mathbin{\mathchoice{\mysetminusD}{\mysetminusT}{\mysetminusS}{\mysetminusSS}}}

% miscellaneous commands
\newcommand{\defn}[1]{{\color{mylightblue}{#1}}}
\newcommand{\toDo}{{\bf\color{red} TODO}}
\newcommand{\toCite}{{\bf\color{green} CITE}}
\newcommand*{\vertbar}{\rule[-1ex]{0.5pt}{2.5ex}} % for matrices with column vectors
\newcommand*{\horzbar}{\rule[.5ex]{2.5ex}{0.5pt}} % for matrices with row vectors
\newcommand{\myblue}[1]{{\color{iceberg}{#1}}}
\newcommand{\myorange}[1]{{\color{burntorange}{#1}}}
\newcommand{\mygreen}[1]{{\color{applegreen}{#1}}}
\newcommand{\myred}[1]{{\color{darkcandyapplered}{#1}}}

% ferrer's diagram
\newcommand{\fdiagram}[1]{
    \begin{tikzpicture}[scale=.7]
        \fill foreach \Z [count=\Y] in {#1}
        {foreach \X in {1,...,\Z} 
        {(\X,-\Y) circle[radius=3pt]}};
    \end{tikzpicture}
}

%
\usepackage{pgfplots}
\pgfplotsset{compat=1.18}

%
\makeatletter
\def\mathcolor#1#{\@mathcolor{#1}}
\def\@mathcolor#1#2#3{%
  \protect\leavevmode
  \begingroup
    \color#1{#2}#3%
  \endgroup
}
\makeatother

% 
\newenvironment{nouppercase}{%
  \let\uppercase\relax%
  \renewcommand{\uppercasenonmath}[1]{}}{}

% titlepage
\title{226: Θεωρία Δακτυλίων και Modules \\ Ασκήσεις με Ενδεικτικές Λύσεις}
\author[Β.~Δ. Μουστακας]{Βασίλης Διονύσης Μουστάκας \\ Πανεπιστήμιο Κρήτης}
\date{29 Απριλίου 2026}
% \urladdr{\href{https://sites.google.com/view/vasmous}{https://sites.google.com/view/vasmous}}

\begin{document}

\begingroup
\def\uppercasenonmath#1{} % this disables uppercase title
\let\MakeUppercase\relax % this disables uppercase authors
\maketitle
\endgroup

\setcounter{exercise}{31}
% \setcounter{theorem}{5}
\setcounter{equation}{9}
\renewcommand{\thesubsection}{\thesection.\arabic{subsection}}

Σε ότι ακολουθεί υποθέτουμε ότι $F$ είναι ένα σώμα.

\begin{exercise}
  \label{ex:ses_vectorSpaces}
  Για κάθε $n \ge 2$, αν 
  \begin{equation} 
    \label{eq:vectorSpaces_ses}
    \begin{tikzcd}
      0 \arrow[r] & V_n \arrow[r, "\phi_n"] & \cdots \arrow[r] & V_1 \arrow[r, "\phi_1"] & V_0 \arrow[r] & 0
    \end{tikzcd}
  \end{equation}
  είναι μια ακριβή ακολουθία διανυσματικών χώρων πεπερασμένης διάστασης πάνω από το $F$, τότε 
  \[
  \sum_{i=0}^n (-1)^i\dim_F(V_i) = 0.
  \]
\end{exercise}

\begin{proof}[Λύση]
  Θα κάνουμε επαγωγή ως προς $n$. Για $n=2$, η ακολουθία \eqref{eq:vectorSpaces_ses} παίρνει την μορφή 
  \[
  \begin{tikzcd}
      0 \arrow[r] & V_2 \arrow[r, "\phi_2"] & V_1 \arrow[r, "\phi_1
    "] & V_0 \arrow[r] & 0.
  \end{tikzcd}
  \]
  Από την παρατήρηση μετά το Θεώρημα 4.12 έπεται ότι 
  \[
  \dim_F\left(
    V_0
  \right) = \dim_F\left(
    V_1
  \right) - \dim_F\left(
    \rmIm(\phi_2)
  \right).
  \]
  Όμως, επειδή η ακολουθία είναι ακριβή στο $V_2$ έπεται ότι η $\phi_2$ είναι μονομορφισμός και γι αυτό $\dim_F(V_2) = \dim_F\left(
    \rmIm(\phi_2)
  \right)$. Επομένως, η παραπάνω σχέση γίνεται 
  \[
  \dim_F\left(
    V_0
  \right) = \dim_F\left(
    V_1
  \right) - \dim_F\left(
    V_2
  \right)
  \] 
  η οποία είναι ισοδύναμη με την ζητούμενη. Αυτή είναι η βάση της επαγωγής.

  Έστω $n > 2$ και υποθέτουμε ότι το ζητούμενο ισχύει για κάθε ακριβή ακολουθία όπως της εκφώνησης με $n$ όρους. Επειδή η ακολουθία \eqref{eq:vectorSpaces_ses} είναι ακριβή στα $V_0$ και $V_1$, έπεται ότι η $\phi_1$ είναι επιμορφισμός και ότι $\rmIm(\phi_2) = \ker(\phi_1)$, αντίστοιχα. Συνεπώς, μπορούμε να την \textquote{χωρίσουμε} σε δύο ακριβής ακολουθίες 
  \[
  \begin{tikzcd}
      0 \arrow[r] & V_n \arrow[r, "\phi_n"] & \cdots \arrow[r, "\phi_3"] & V_2 \arrow[r, two heads] & \ker(\phi_1) \arrow[r] & 0
    \end{tikzcd}
  \]
  και 
  \[
  \begin{tikzcd}
      0 \arrow[r] & \ker(\phi_1) \arrow[r, hook] & V_1 \arrow[r, "\phi_1
    "] & V_0 \arrow[r] & 0
  \end{tikzcd}
  \]
  (γιατί;). Από την επαγωγική υπόθεση έπεται ότι 
  \begin{align*}
    \dim_F\left(
      \ker(\phi_1)
    \right) - \sum_{i=2}^n (-1)^i \dim_F\left(
      V_i
    \right) &= 0 \\ 
    \dim_F\left(
      V_0
    \right) - \dim_F\left(
      V_1
    \right) + \dim_F\left(
      \ker(\phi_1)
    \right) &= 0.
  \end{align*}
  Η ζητούμενη ισότητα προκύπτει αφαιρώντας τις δυο ισότητες κατά μέλη.
\end{proof}

\begin{remark}
  Η λύση της Άσκησης \ref{ex:ses_vectorSpaces} συνεχίζει να ισχύει και με την (γενίκοτερη) υπόθεση κάθε $V_i$ να είναι ελεύθερο πρότυπο πεπερασμένης τάξης πάνω από μια περιοχή κύριων ιδεωδών (γιατί;).
\end{remark}

Σε ότι ακολουθεί, υποθέτουμε ότι κάθε διαβαθμισμένη $F$-άλγεβρα έχει πεπερασμένης διάστασης ομογενείς συνιστώσες. Δηλαδή, αν $A$ είναι μια $F$-άλγεβρα με διαβάθμιση
\[
A = A_0 \oplus A_1 \oplus \cdots, 
\]
τότε $\dim_F(A_d) < \infty$, για κάθε $d \ge 0$.

\begin{exercise}
  \label{ex:graded_ideal}
  Έστω $A$ μια διαβαθμισμένη $F$-άλγεβρα και $I$ ένα ιδεώδες. Το $I$ ονομάζεται \defn{ομογενές} (ή \defn{διαβαθμισμένο}) αν 
  \[
  I = \left(
    I \cap A_0
  \right) \oplus \left(
    I \cap A_1
  \right) \oplus \cdots.
  \]
  Να δείξετε ότι αν το $I$ είναι ομογενές, τότε το $A/I$ είναι διαβαθμισμένη $F$-άλγεβρα με διάβάθμιση
  \[
  A/I \cong A_0/I_0 \oplus A_1/I_1 \oplus \cdots.
  \]
\end{exercise}

\begin{proof}[Yπόδειξη]
  Η απεικόνιση 
  \begin{align*}
    A_0 \oplus A_1 \oplus \cdots &\to A_0/I_0 \oplus A_1/I_1 \oplus \cdots \\ 
  (a_0, a_1, \dots) &\mapsto (a_0 + I_0, a_1 + I_1, \dots)
  \end{align*}
  είναι επιμορφισμός $A$-αλγεβρών με πυρήνα το $I$.
\end{proof}

\begin{exercise}
  \label{ex:graded_ideal_example}
  Έστω $A = F[x,y,z]$ και $I = (xy)$. Να δείξετε ότι το $I$ είναι ομογενές και να υπολογίσετε την αντίστοιχη σειρά Hilbert.
\end{exercise}

\begin{proof}[Λύση]
  Από την Άσκηση 31 γνωρίζουμε ότι 
  \[
  A = F\{1\} \oplus F\{x,y,z\} \oplus F\{x^2,y^2,z^2,xy,xz,yz\} \oplus \cdots.
  \]
  Συνεπώς, 
  \[
  I = 0 \oplus 0 \oplus F\{\mathcolor{burntorange}{xy}\} \oplus F\{\underbrace{x\mathcolor{burntorange}{xy}}_{= \ x^2y}, \underbrace{\mathcolor{burntorange}{xy}y}_{= \ xy^2}, \mathcolor{burntorange}{xy}z\} \oplus \cdots
  \]
  και γι αυτό 
  \begin{align*}
  \Hilb(I;t) &= \sum_{d \ge 0} \dim_F(A) t^{d+2} \\
  &= t^2 \sum_{d \ge 0} \dim_F(A) t^{d} \\
  &= \frac{t^2}{(1-t)^2},
  \end{align*}
  όπου η τελευταία ισότητα έπεται από την Άσκηση 31.
\end{proof}

\begin{remark}
  Το ιδεώδες $I$ της Άσκησης \ref{ex:graded_ideal_example} είναι ελεύθερο πάνω από το $A$ (γιατί;). Ουσιαστικά, πρόκειται για ένα αντίγραφο της $A$ \textquote{μετακινημένο} κατά δύο θέσεις προς τα δεξιά.
\end{remark}

\begin{exercise}
  \label{ex:free_module_generated_in_degree}
  Έστω $A$ μια διαβαθμισμένη $F$-άλγεβρα με διάσπαση 
  \[
  A = A_0 \oplus A_1 \oplus \cdots.
  \] 
  Αν $f \in A_k$ ένα \emph{ομογενές στοιχείο} βαθμού $k$, τότε να δείξετε ότι
  \[
  \Hilb((f);t) = t^k \, \Hilb(A;t).
  \]
\end{exercise}

\begin{proof}[Λύση]
  Έστω $I = (f)$. Επειδή το $I$ είναι ιδεώδες και η $A$ είναι διαβαθμισμένη έπεται ότι αν $a \in A_d$, τότε $af \in I \cap A_{d+k}$, για κάθε $d \ge 0$. Συνεπώς, 
  \[
  I \cap A_d = 
  \begin{cases}
    0, &\ \text{αν $d < k$} \\
    A_{d-k}\{f\}, &\ \text{αν $d \ge k$}
  \end{cases}
  \]
  και γι αυτό 
  \[
  \dim_F\left(
    I \cap A_d
  \right) = \dim_F\left(
    A_{d-k}
  \right)
  \]
  για κάθε $d \ge k$. Άρα, 
  \begin{align*}
  \Hilb(I;t) &= \sum_{d \ge k} \dim_F\left(
    A_{d-k}
  \right) t^d \\ 
  &= \sum_{d\ge0} \dim_F\left(
    A_{d}
  \right) t^{d+k} \\ 
  &= t^k \, \Hilb(A;t).
  \end{align*}
\end{proof}

\begin{remark}
  Το ιδεώδες $(f)$ είναι ένα ελεύθερο $A$-πρότυπο τάξης $1$, δηλαδή ένα (ισομορφικό) αντίγραφο του $A$. Επειδή ο γεννήτοράς του έχει βαθμό $k$, η Άσκηση \ref{ex:free_module_generated_in_degree} μας πληροφορεί ότι η διαβάθμιση είναι ουσιαστικά 
  \[
  \underbrace{0 \oplus 0 \oplus \cdots \oplus 0}_{\text{$k$ φορές}} \oplus A_0 \oplus A_1 \oplus \cdots.
  \] 
  Αυτό συνήθως συμβολίζεται ως $A(-k)$.
\end{remark}

\begin{exercise}
  \label{ex:hilbertSeries_xy}
  Έστω $A = F[x,y,z]$ και $I = (xy)$. Να υπολογιστεί η σειρά Hilbert του $A/I$.
\end{exercise}

\begin{proof}[Πρώτη Λύση]
  Από την Άσκηση \ref{ex:free_module_generated_in_degree} γνωρίζουμε ότι 
  \[
  I \cong 0 \oplus 0 \oplus A_0 \oplus A_1 \oplus \cdots
  \]
  και γι αυτό, από την Άσκηση \ref{ex:graded_ideal}, έπεται ότι 
  \[
  A/I = \bigoplus_{d \ge 0} \left(
    A_d/I_d
  \right) 
  \cong A_0 \oplus A_1 \oplus \frac{A_2}{A_0} \oplus \frac{A_3}{A_1} \oplus \cdots.
  \]
  Άρα,
  \begin{align*}
  \Hilb(A/I;t) &= \sum_{d \ge 0} \dim_F(A_d) t^d - \sum_{d\ge2}\dim_F(A_{d-2}) t^d \\
  &= \sum_{d \ge 0} \dim_F(A_d) t^d - t^2\sum_{d\ge0}\dim_F(A_{d}) t^d \\
  &= \frac{1-t^2}{(1-t)^2} \\
  &= \frac{1+1}{1-t}.
  \end{align*}
  (Πώς δικαιολογείται η τελευταία ισότητα στον δακτύλιο των τυπικών δυναμοσειρών $\ZZ[\![t]\!]$;)
\end{proof}

\begin{proof}[Δεύτερη Λύση]
  Θεωρούμε την βραχεία ακριβή ακολουθία 
  \[
  \begin{tikzcd}
      0 \arrow[r] & I \arrow[r, hook] & A \arrow[r, two heads] & A/I \arrow[r] & 0
    \end{tikzcd}.
  \]
  Επειδή το $I$ είναι ομογενές ιδεώδες και κατά συνέπεια το ίδιο ισχύει και για το $A/I$ (από την Άσκηση \ref{ex:graded_ideal}), αυτή επάγει ακριβείς ακολουθίες 
  \[
  \begin{tikzcd}
      0 \arrow[r] & I_d \arrow[r, hook] & A_d \arrow[r, two heads] & A_d/I_d\arrow[r] & 0
    \end{tikzcd}
  \]
  για κάθε $d \ge 0$. Από την Άσκηση \ref{ex:ses_vectorSpaces}, έπεται ότι 
  \[
  \Hilb(A/I;t) = \Hilb(A;t) - \Hilb(I;t) = \frac{1-t^2}{(1-t)^2} =  \frac{1+1}{1-t}.
  \]
\end{proof}

\begin{exercise}
  \label{ex:hilbertSeries_of_quotient_rings}
  Έστω $A$ μια $F$-άλγεβρα και $I$ ένα ομογενές ιδεώδες. Να δείξετε ότι 
  \[
  \Hilb(A/I;t) = \Hilb(A;t) - \Hilb(I;t).
  \]
\end{exercise}

\begin{proof}[Υπόδειξη]
  Μιμηθείτε την δεύτερη λύση της Άσκησης \ref{ex:hilbertSeries_xy}.
\end{proof}

\begin{remark}
  Η Άσκηση \ref{ex:hilbertSeries_of_quotient_rings} ανάγει τον υπολογισμό της σειράς Hilbert μιας $F$-άλγεβρας $A$ στον υπολογισμό των αντίστοιχων σειρών Hilbert για τα $I$ και $A/I$, για \emph{οποιοδήποτε} ιδεώδες $I$.
\end{remark}

\end{document}